\section{Pilot Study II: Bias-Aware Optimization}
\label{sec:optimizer_pilot}

Adaptive optimizers such as Adam and Muon maintain
second-moment statistics that are highly sensitive to extreme gradients.
Section~\ref{subsec:mean_second_moment} showed that
mean bias dominates second-moment tails.
This motivates bias-aware optimization.

\subsection{Hypothesis}

\begin{quote}
\emph{If mean bias inflates second-moment statistics,
then removing activation mean should reduce effective learning-rate suppression
and partially substitute spectral optimizers such as Muon and Spectra.}
\end{quote}

\subsection{Method}

We evaluate three variants:

\begin{itemize}
\item Standard AdamW
\item Spectra (spike-aware optimizer)
\item AdamW + source-level mean removal
\end{itemize}

Mean removal is applied only at embedding activation.

\subsection{Experimental Setup}

\begin{itemize}
\item Model: 1B parameters.
\item Training: 100B tokens.
\item Full precision (BF16).
\item Metrics:
  \begin{itemize}
  \item Training loss curve
  \item Gradient variance
  \item Effective learning-rate distribution
  \end{itemize}
\end{itemize}

\subsection{Results}

\begin{itemize}
\item Mean removal reduces $p99.9$ of second moment $v$ by $4$--$10\times$.
\item Effective learning-rate distribution becomeFs more uniform.
\item Convergence speed improves relative to baseline AdamW.
\item Performance approaches that of Spectra without per-layer SVD.
\end{itemize}

\subsection{Interpretation}

Second-moment inflation is primarily driven by coherent mean drift.
Source-level bias conditioning reduces this inflation
without explicit spectral decomposition.
This suggests that spike-aware optimizers may be approximated
by simpler bias-aware preprocessing.

\subsection{Limitation}

Unlike Spectra, mean removal does not adaptively track
rotating spike subspaces.
Thus in later training phases,
residual low-rank structure may still influence optimization.